% Focused insertion snippets for Totani_Halo_v2_copyedited.tex (October 7, 2026).
% Copy each block at its own anchor; do not input this complete file at once.
% Use only production PDFs from a scientifically validated compact bundle.
% Fixture PDFs visibly marked TEST DATA are exclusively layout proofs.
% Keep the existing title, abstract, prose, equations, figures, tables and
% refs.bib unchanged. All citation keys below already exist in refs.bib.
% The current editorial REVTeX text width is 510 TeX pt = 7.05687 inches.
% Put these production PDFs in the manuscript's figures/ directory, which is
% already on its graphicspath. Review the generated bundle notes before use.

% BEGIN TEMPLATE OVERVIEW INSERTION
% Anchor: end of Section II, after subsection NFW template (sec:halo),
% immediately before \section{Statistical analysis} (sec:method).
Figure~\ref{fig:template_overview} shows the NFW and bubble templates
together with the masks that define the fitted region.
\begin{figure*}[t]
\centering
\includegraphics[width=\textwidth]{template_overview.pdf}
\caption{Shown here are the spatial templates whose fitted amplitudes are
compared in this work, on the $120^\circ\times120^\circ$ analysis grid
before exposure weighting and PSF convolution: (a) the NFW annihilation
template, $J(\ell,b)$ of Eq.~(\ref{eq:jmap}); (b) the flat \fermi\ bubbles
template \cite{Fermi-LAT:2014sfa}; (c) the structured \fermi\ bubbles
template \cite{Fermi-LAT:2017opo}; and (d) the
X-ray-brightness-weighted eROSITA bubbles template
\cite{Predehl:2020kyq,Zhang:2024mcd}. The NFW template overlaps the
\fermi\ bubbles over much of the fitted region at $|b|\lesssim30^\circ$,
and the eastern ($\ell>0$) half of the eROSITA bubbles template, which
reduces the gamma-ray asymmetry in Sec.~\ref{sec:erositafreedom}, is a
reconstruction rather than native survey data. Each map is divided by
its maximum over the full template, including the masked plane, so that
the NFW template within the fitted region reaches only $5.3\times10^{-3}$
of its central value; the color scales of (a) and (d) are stretched for
display, and the northern and southern halves of the eROSITA template
retain their independent normalizations. The solid (black) line is the
boundary of the flat \fermi\ bubbles, the dashed (orange) line in (a)
encloses the region outside the bubbles used in Sec.~\ref{sec:outside},
and the dashed (gray) line in (d) separates the native western from the
reconstructed eastern eROSITA data. Light gray marks the excluded plane,
$|b|<10^\circ$, and the 4FGL-DR4 extended sources masked within twice
their semi-major axes \cite{Ballet:2023qzs}. The dark disk at
$(\ell,b)\approx(-1^\circ,+24^\circ)$ in (d) is the region around
Sco~X-1 outside the eROSITA coverage, where the template weight is zero.}
\label{fig:template_overview}
\end{figure*}
% END TEMPLATE OVERVIEW INSERTION

% BEGIN REGIONAL NORMALIZATIONS INSERTION
% Anchor: Section V, subsection Quadrant likelihood tests and simulated halos
% (sec:quadrantfinal), after the paragraph ending "a more extreme significance."
% and before the paragraph beginning "Using twelve brightness-weighted ...".
% Preserve the existing prose's model-set p=0.5--0.9 range as model-set prose;
% the plot's baseline panels must use only their own exported measurements.
Figure~\ref{fig:regional_nfw_normalizations} shows the regional NFW
coefficients that underlie these likelihood comparisons.
\begin{figure*}[t]
\centering
\includegraphics[width=\textwidth]{regional_nfw_normalizations.pdf}
\caption{Shown here are the NFW coefficients fitted independently in four
sky regions at 20.8\,GeV for Totani's baseline diffuse model in the
single-gas decomposition, with $2^\circ$ cells. The left column uses the
standard quadrants and the right column the partition rotated by
$45^\circ$, both drawn at top with the excluded plane in white; E and W
denote $\ell>0$ and $\ell<0$. From top to bottom, the rows use the
structured \fermi\ bubbles template with uniform northern and southern
eROSITA bubbles templates, the flat \fermi\ bubbles template with the
same eROSITA templates, and the flat \fermi\ bubbles template with twelve
X-ray-brightness-weighted eROSITA bubbles pieces. Only with the twelve
pieces do the standard quadrants agree with a common normalization,
whereas the rotated partition disfavors a common normalization in every
row, with the east and west wedges near the plane containing little
NFW-correlated emission (Sec.~\ref{sec:quadrantfinal}). The filled
(black) circles are the regional coefficients with 68\%
profile-likelihood intervals, in which the other regional and all
foreground amplitudes are refitted; the solid (blue) line is the common
coefficient fitted to the same data; and the dotted (gray) line marks
zero. The regional templates are cut from the NFW intensity map before
PSF convolution, so that equal coefficients reproduce the full folded
template. The coefficient multiplies $J$ of Eq.~(\ref{eq:jmap}) to give
the differential intensity for a flat spectrum across the bin, in units
of $10^{-35}\,\mathrm{cm^3\,s^{-1}\,GeV^{-2}\,MeV^{-1}\,sr^{-1}}$. In
each panel, $Q_6$ is the sum of $\dl$ between the common and regional
fits over the six bins from 5.6 to 130\,GeV, and the count is the number
of 200 null simulations of that fixed partition that exceed it. These
simulations were drawn from a reference foreground fit with a different
GALPROP model from the baseline shown, and therefore calibrate this
configuration only approximately.}
\label{fig:regional_nfw_normalizations}
\end{figure*}
% END REGIONAL NORMALIZATIONS INSERTION
